\section{Canonical completeness for serial, symmetric and transitive systems}

\begin{annexitem}{annex:p8.mce.1}{thm:p08-mce-soundness}{Soundness, exact frame correspondences and consistency}
\begin{proof}
A propositional tautology remains true when its propositional variables
are replaced by modal formulas: at a fixed world each substituted formula
has a Boolean truth value. For K, if every successor satisfies \(A\to B\)
and every successor satisfies \(A\), every successor satisfies \(B\).
Thus the distribution axiom is true on every frame.

On a serial frame, if \(\Box A\) is true at \(w\), take a successor
\(v\), whose existence seriality asserts. It satisfies \(A\), so
\(\Diamond A\) is true at \(w\). On a reflexive frame \(wRw\)
makes \(\Box A\to A\) true. On a symmetric frame, suppose \(A\)
is true at \(w\). For any \(v\) with \(wRv\), symmetry gives
\(vRw\), so \(\Diamond A\) is true at \(v\). Hence
\(A\to\Box\Diamond A\) is true at \(w\). On a transitive frame,
suppose \(\Box A\) is true at \(w\) and \(wRv\). If \(vRu\),
then \(wRu\), so \(A\) is true at \(u\). Thus \(\Box A\) is
true at every successor \(v\), proving \(\Box A\to\Box\Box A\).

For the converse frame implications, fix an atom \(p\). If \(w\)
has no successor, then \(\Box\topv\) is true and
\(\Diamond\topv\) false there, so D fails. If \(wRw\) fails,
put \(V(p)=W\setminus\{w\}\). Every successor satisfies \(p\)
while \(w\) does not, refuting T. If \(wRv\) but not \(vRw\),
put \(V(p)=\{w\}\). Then \(p\) is true at \(w\), but
\(\Diamond p\) is false at its successor \(v\), refuting B.
If \(wRv\) and \(vRu\) but not \(wRu\), put
\(V(p)=\{x:wRx\}\). Then \(\Box p\) is true at \(w\), while
\(p\) is false at \(u\), so \(\Box p\) is false at \(v\) and
\(\Box\Box p\) is false at \(w\). This refutes 4. These valuations
prove each converse under the all-valuations definition of frame validity.

For no-premise derivations, induct on the finite proof together with the
separate proofs attached to necessitation steps. Modus ponens preserves
truth. At a necessitation step, the induction hypothesis makes its premise
theorem valid at every world on every frame in the class. It is therefore
true at every successor of the current world, and its box is valid.
All the axioms of the three selected systems have just been checked on
their respective frame classes. Hence every no-premise theorem is valid.
Now fix a structure, a world satisfying \(\Gamma\), and a finite
derivation from \(\Gamma\). Premise lines are true there, axiom lines
are true by the preceding checks, and modus ponens preserves truth there.
A necessitated line is true there by the already proved no-premise result;
its separate theorem proof uses no assumption from \(\Gamma\).
Induction on these lines proves local soundness for arbitrary \(\Gamma\).

The one-world frame with its single reflexive edge is serial, reflexive,
symmetric and transitive. On any valuation over this frame \(\bottom\)
is false. Soundness therefore excludes a no-premise derivation of
\(\bottom\) in each of \(KD,KTB,K4\), proving their consistency.
\end{proof}
\end{annexitem}

\begin{annexitem}{annex:p8.mce.2}{thm:p08-mce-deduction}{Finite support, the deduction theorem and boxed finite consequences}
\begin{proof}
A finite derivation has only finitely many premise lines. Its separate
theorem proofs contribute no premises. Keeping those finitely many premise
formulas gives the same derivation, which proves finite support.

For the forward deduction implication, induct on a derivation from
\(\Gamma\cup\{A\}\). For a line equal to \(A\), the required
\(A\to A\) is a tautology. For a line \(C\) that is an axiom or
a member of \(\Gamma\), derive \(A\to C\) using \(C\) and the
tautology \(C\to(A\to C)\). For a modus-ponens line \(D\)
obtained from \(C\) and \(C\to D\), the induction hypothesis gives
\(A\to C\) and \(A\to(C\to D)\). The tautology
\[
 (A\to(C\to D))\to((A\to C)\to(A\to D))
\]
and two applications of modus ponens give \(A\to D\). Finally, if a
line is \(\Box C\), its separate no-premise proof of \(C\) also
licenses \(\Box C\) in the derivation from \(\Gamma\). The
tautology \(\Box C\to(A\to\Box C)\) gives the desired line.
There is no case in which a premise-dependent proof of \(C\) is boxed.
The reverse deduction implication is modus ponens after adding \(A\)
to the premises.

If \(\Gamma\cup\{\neg A\}\) were inconsistent, deduction would
give \(\Gamma\derives_L\neg A\to\bottom\). The formula
\((\neg A\to\bottom)\to A\) is a classical tautology, so
\(\Gamma\derives_L A\), contradicting the hypothesis. More generally,
if \(X\) is consistent, at least one of \(X\cup\{A\}\) and
\(X\cup\{\neg A\}\) is consistent: inconsistency of both gives,
by deduction and classical tautologies, both \(\neg A\) and \(A\)
from \(X\), and hence \(\bottom\).

For the boxed claim, repeated deduction transforms a derivation of \(B\)
from \(C_1,\ldots,C_n\) into a no-premise derivation of
\(C_1\to(C_2\to\cdots(C_n\to B)\cdots)\). Necessitate this
theorem once. Apply the K schema to its first implication, and modus
ponens with the premise \(\Box C_1\), to obtain the box of the
remaining implication. Repeating with \(\Box C_2,\ldots,\Box C_n\)
gives \(\Box B\). When \(n=0\), the original derivation is already
a no-premise proof of \(B\), so a single theorem-only necessitation
suffices. This also covers repetitions or an empty selection of boxed
antecedents.
\end{proof}
\end{annexitem}

\begin{annexitem}{annex:p8.mce.3}{thm:p08-mce-lindenbaum}{Extension of arbitrary consistent theories and all Boolean cases}
\begin{proof}
Fix an enumeration \(A_0,A_1,\ldots\) of all finite modal formulas.
For a consistent set \(X\), define \(X_0=X\). Given \(X_n\), set
\(X_{n+1}=X_n\cup\{A_n\}\) if that set is consistent, and set
\(X_{n+1}=X_n\cup\{\neg A_n\}\) otherwise. The consistency
alternative proved in Annex~\ref{annex:p8.mce.2} shows by induction that
each stage is consistent. Every stage contains the preceding one.
This is a uniquely specified classical set recursion; no decision algorithm
for the consistency test and no choice principle are asserted.

Put \(w=\bigcup_{n\in\N}X_n\). If \(w\) were inconsistent,
finite support would give a finite set of premise formulas witnessing it.
Each such formula belongs to some stage. Choose for each its least stage
index and take the maximum of this finite set of natural numbers; if the
premise set is empty take index zero. All the premises then lie in one
consistent stage, a contradiction. For every formula \(A\), the stage
at which it is enumerated puts either \(A\) or \(\neg A\) into
\(w\). They cannot both belong, since their joint presence would derive
\(\bottom\). If \(A\notin w\), the presence of \(\neg A\)
makes every extension containing \(A\) inconsistent. Therefore \(w\)
is maximal and extends \(X\), regardless of whether \(X\) was finite.

Now let \(w\) be any maximal consistent set. If \(A\notin w\),
maximality makes \(w\cup\{A\}\) inconsistent. Deduction gives
\(w\derives_L A\to\bottom\), hence
\(w\derives_L\neg A\). To establish closure without presupposing it,
suppose \(w\derives_L C\) and \(C\notin w\). A derivation of
\(\bottom\) from \(w\cup\{C\}\) exists by maximality.
Replace each use of premise \(C\) by the finite derivation from \(w\).
The result is a derivation of \(\bottom\) from \(w\), impossible.
Thus \(w\derives_L C\) implies \(C\in w\); the converse is the
one-line premise derivation. In particular, all no-premise theorems belong
to \(w\), and it is closed under modus ponens. The above derivation of
\(\neg A\) now places \(\neg A\) in \(w\). Consistency prevents
both \(A\) and \(\neg A\), proving the negation equivalence.
It also excludes \(\bottom\), while the tautology \(\neg\bottom\)
puts \(\topv\) in \(w\).

If \(A\land B\in w\), its tautological projections and closure give
\(A,B\in w\). Conversely, the tautology
\(A\to(B\to A\land B)\) combines their membership.
If either disjunct belongs, the corresponding tautological introduction
puts \(A\lor B\) in \(w\). If neither belongs, then
\(\neg A,\neg B\in w\); together with \(A\lor B\) these derive
\(\bottom\). Thus membership of the disjunction forces at least one
disjunct. If \(A\to B\in w\) and \(A\in w\), closure gives
\(B\in w\). Conversely, if \(A\notin w\), use
\(\neg A\in w\) and the tautology \(\neg A\to(A\to B)\).
If \(B\in w\), use \(B\to(A\to B)\). These are exactly the
two alternatives for implication in the statement, completing every
primitive Boolean membership case.
\end{proof}
\end{annexitem}

\begin{annexitem}{annex:p8.mce.4}{thm:p08-mce-existence}{Successors witnessing failure of a box}
\begin{proof}
Suppose \(\neg\Box A\in w\). If
\(U_L(w)\cup\{\neg A\}\) were inconsistent, deduction and the
classical double-negation tautology would give
\(U_L(w)\derives_L A\). By finite support there would be
\(C_1,\ldots,C_n\in U_L(w)\), possibly no such formulas, from
which \(A\) is derivable. The boxed-consequence part of
Lemma~\ref{thm:p08-mce-deduction} then gives \(\Box A\) from
\(\Box C_1,\ldots,\Box C_n\). Each boxed premise belongs to
\(w\) by definition of \(U_L(w)\), so closure puts \(\Box A\)
in \(w\), contradicting its consistency. The set is therefore
consistent. Lindenbaum extension gives \(v\) containing it, and hence
\(wR_Lv\) and \(\neg A\in v\).

If \(\Diamond B\in w\), then \(\neg\Box\neg B\in w\).
The result just proved, with \(A=\neg B\), gives a successor containing
\(\neg\neg B\), and closure under classical double-negation
elimination puts \(B\) there. Conversely, suppose \(wR_Lv\) and
\(B\in v\). If \(\Diamond B\notin w\), Boolean maximality and
its defining abbreviation put \(\Box\neg B\) in \(w\).
Accessibility then puts \(\neg B\) in \(v\), contradicting
\(B\in v\). Thus \(\Diamond B\in w\).
\end{proof}
\end{annexitem}

\begin{annexitem}{annex:p8.mce.5}{thm:p08-mce-truth}{Full structural induction for canonical truth}
\begin{proof}
Induct on the finite construction of \(A\), with the assertion made
simultaneously for all worlds. For an atom \(p\), the equivalence is
the definition of \(V_L(p)\). For \(\bottom\), both truth and
membership are false. If the assertion holds for \(B\), then
\(\neg B\) is true at \(w\) exactly when \(B\) is not true
there, exactly when \(B\notin w\), exactly when \(\neg B\in w\).

For \(B\land C\), the semantic requirement is truth of both
subformulas, which by induction is membership of both and by
Lemma~\ref{thm:p08-mce-lindenbaum} is membership of the conjunction.
For \(B\lor C\), the semantic requirement is truth of at least one
subformula, which is membership of at least one and therefore membership
of the disjunction. For \(B\to C\), the semantic requirement is
failure of \(B\) or truth of \(C\), which is \(B\notin w\) or
\(C\in w\), exactly membership of the implication. This treats each
primitive binary Boolean connective separately.

For \(\Box B\), if \(\Box B\in w\), every successor \(v\)
contains \(B\) by accessibility, so every successor satisfies \(B\)
by induction. Hence the box is true. If \(\Box B\notin w\), then
\(\neg\Box B\in w\), and Lemma~\ref{thm:p08-mce-existence}
gives a successor \(v\) with \(\neg B\in v\), hence \(B\notin v\).
The induction hypothesis makes \(B\) false at that successor, so the
box is false. There are no other primitive cases. Expanding
\(\Diamond B=\neg\Box\neg B\) gives its truth equivalence; the
existential form is also exactly the second assertion of the existence
lemma combined with induction.
\end{proof}
\end{annexitem}

\begin{annexitem}{annex:p8.mce.6}{thm:p08-mce-canonical-frames}{The D, T, B and 4 canonical-frame arguments}
\begin{proof}
For \(KD\), \(\topv\) is a theorem, so theorem-only
necessitation makes \(\Box\topv\) a theorem. Its D instance gives
\(\Diamond\topv\). Both theoremhood and maximal closure therefore
put \(\Diamond\topv\) into every world \(w\). The existence
lemma supplies a successor \(v\), proving seriality. This is an
existence assertion at each world; no simultaneous choice of successors
is needed to define \(R_{KD}\).

For \(KTB\), if \(\Box A\in w\), its T instance and modus
ponens give \(A\in w\). Thus \(U_{KTB}(w)\subseteq w\),
so \(wR_{KTB}w\). For symmetry suppose \(wR_{KTB}v\), and let
\(\Box A\in v\). If \(A\notin w\), then \(\neg A\in w\).
The B instance with \(\neg A\) gives \(\Box\Diamond\neg A\in w\),
so accessibility gives \(\Diamond\neg A\in v\). Its expansion is
\(\neg\Box\neg\neg A\). On the other hand the tautological
\(A\to\neg\neg A\), followed by theorem-only necessitation and
K, yields \(\Box A\to\Box\neg\neg A\). The assumption
\(\Box A\in v\) puts \(\Box\neg\neg A\) in \(v\), a
contradiction. Therefore \(A\in w\). This proves
\(U_{KTB}(v)\subseteq w\), or \(vR_{KTB}w\). Notice that the
symmetry argument uses B, while reflexivity used the independent T
component; they are not merged into one axiom.

For \(K4\), suppose \(wR_{K4}v\) and \(vR_{K4}u\). For any
\(\Box A\in w\), its 4 instance puts \(\Box\Box A\in w\).
The first edge then puts \(\Box A\in v\), and the second puts
\(A\in u\). Thus \(U_{K4}(w)\subseteq u\), proving
\(wR_{K4}u\). No T or D axiom was used. Combined with nonemptiness
from Definition~\ref{def:p08-mce-canonical-model}, these arguments put
each canonical frame in its stated class.
\end{proof}
\end{annexitem}

\begin{annexitem}{annex:p8.mce.7}{thm:p08-kd-strong-completeness}{Strong local adequacy on serial frames}
\begin{proof}
The implication from derivability to local consequence is
Lemma~\ref{thm:p08-mce-soundness}. Conversely, if
\(\Gamma\not\derives_{KD}A\),
Lemma~\ref{thm:p08-mce-deduction} makes
\(\Gamma\cup\{\neg A\}\) consistent. Extend that entire set to
a maximal \(KD\)-consistent world \(w\). The truth lemma makes
every formula in \(\Gamma\) true at \(w\) and \(A\) false.
The canonical frame is nonempty and serial by
Lemma~\ref{thm:p08-mce-canonical-frames}. This is a countermodel to
the proposed local consequence on exactly the required frame class.
For satisfiability of a consistent \(\Gamma\), extend \(\Gamma\)
itself instead and apply the same truth and frame lemmas. At no point is
\(\Gamma\) replaced by a finite subset in the model construction.
\end{proof}
\end{annexitem}

\begin{annexitem}{annex:p8.mce.8}{thm:p08-ktb-strong-completeness}{Strong local adequacy on reflexive symmetric frames}
\begin{proof}
Soundness holds by Lemma~\ref{thm:p08-mce-soundness}, using both the
T and B checks. If \(\Gamma\not\derives_{KTB}A\), extend the
consistent set \(\Gamma\cup\{\neg A\}\) to a maximal world
\(w\). In the \(KTB\) canonical structure all members of
\(\Gamma\) are true at \(w\), while \(A\) is false, by the
truth lemma. The frame is reflexive by T and symmetric by the separate
B argument in Lemma~\ref{thm:p08-mce-canonical-frames}; it is nonempty
because it contains \(w\). It therefore witnesses failure of local
consequence in \(\mathcal C_{KTB}\), proving the reverse implication.
A maximal extension of any consistent \(\Gamma\) supplies its
satisfying world in the same class. None of these steps asserts that
the canonical relation is transitive or changes \(KTB\) to \(KB\).
\end{proof}
\end{annexitem}

\begin{annexitem}{annex:p8.mce.9}{thm:p08-k4-strong-completeness}{Strong local adequacy on transitive frames}
\begin{proof}
Use Lemma~\ref{thm:p08-mce-soundness} for soundness. If
\(\Gamma\not\derives_{K4}A\), extend
\(\Gamma\cup\{\neg A\}\) to a maximal \(K4\)-consistent world
\(w\). The truth lemma makes \(w\) satisfy every premise and
falsify \(A\). The canonical relation is transitive by the 4 argument
of Lemma~\ref{thm:p08-mce-canonical-frames}, so this countermodel belongs
to \(\mathcal C_{K4}\). This proves completeness for the original
arbitrary premise set. Extending a consistent \(\Gamma\) alone proves
its satisfiability. The model's world set is nonempty; individual worlds
are allowed to have no successors, since neither seriality nor reflexivity
belongs to this frame requirement.
\end{proof}
\end{annexitem}

\begin{annexitem}{annex:p8.mce.10}{thm:p08-mce-separations}{Complete finite separating models}
\begin{proof}
In each construction, exactly the listed edges are present. Atoms other
than \(p\) may all be assigned the empty set.

For \(\not\derives_{KD}\Box p\to p\), take worlds \(w,v\), edges
\((w,v),(v,v)\), and \(V(p)=\{v\}\). Both worlds have a successor,
so the frame is serial. At \(w\), \(\Box p\) is true because its
sole successor is \(v\), but \(p\) is false. KD soundness proves
nonderivability. This frame is also transitive, so it also witnesses that
transitivity does not force reflexivity.

For \(\not\derives_{K4}\Box p\to\Diamond p\), take the single
world \(w\) with no edges. Transitivity holds vacuously. Both
\(\Box p\) and \(\Box\neg p\) are true, while \(\Diamond p\)
is false. K4 soundness proves the claim and separates transitivity from
seriality.

For \(\not\derives_{KTB}\Box p\to\Box\Box p\), take three worlds
\(a,b,c\), all three reflexive edges, and the four edges
\((a,b),(b,a),(b,c),(c,b)\). Put \(V(p)=\{a,b\}\). The frame
is reflexive and symmetric. The successors of \(a\) are exactly
\(a,b\), both satisfying \(p\), so \(\Box p\) is true at
\(a\). But \(bRc\) and \(p\) is false at \(c\), so
\(\Box p\) is false at \(b\). Since \(aRb\),
\(\Box\Box p\) is false at \(a\). Soundness gives the claimed
failure and shows that reflexivity plus symmetry does not impose
transitivity.

For \(\not\derives_{K4}p\to\Box\Diamond p\), take worlds
\(a,b\), edges \((a,a),(a,b),(b,b)\), and \(V(p)=\{a\}\).
The relation is transitive. At \(a\), \(p\) is true, but its
successor \(b\) sees only \(b\), where \(p\) is false.
Thus \(\Diamond p\) is false at \(b\), and
\(\Box\Diamond p\) is false at \(a\). This proves the claim.

To distinguish the symmetry-only B extension from \(KTB\), take
worlds \(a,b\) and edges \((a,b),(b,a)\), with
\(V(p)=\{b\}\). The frame is symmetric, so every B instance is
valid there by the frame argument in Annex~\ref{annex:p8.mce.1}; K
and its rules are sound there by the same argument. At \(a\),
\(\Box p\) is true and \(p\) is false, so T is not derivable
in \(KB\). This supplies a witness for the distinct convention rather
than substituting that system into the \(KTB\) theorem.

Finally, take two worlds \(a,b\), let the relation be the full
\(\{a,b\}^2\), and put \(V(p)=\{a\}\). This frame is serial,
reflexive, symmetric and transitive, hence belongs to all three classes.
At \(a\), the premise \(p\) is true, but the accessible world
\(b\) makes \(\Box p\) false. Local soundness gives
\(\{p\}\not\derives_L\Box p\) for each selected \(L\).
If instead \(p\) is assumed true at every world of a structure, every
successor of every world satisfies \(p\), so \(\Box p\) is true
everywhere. That proves the stated global consequence and exhibits the
precise failure of premise-dependent necessitation for local consequence.
\end{proof}
\end{annexitem}

\begin{annexitem}{annex:p8.mce.11}{thm:p08-mce-canonicity-compactness}{Arbitrary valuations and finite-support consequences}
\begin{proof}
For each selected \(L\), its canonical frame belongs to
\(\mathcal C_L\). Given any valuation on that frame, soundness makes
each \(L\)-theorem true at every world. This proves canonical-frame
validity under every valuation, which is the defined canonicity property.

Suppose \(\Gamma\entails^{\ell}_{\mathcal C_L}A\). The corresponding
strong completeness theorem gives a derivation of \(A\) from
\(\Gamma\). Finite support gives a finite
\(\Gamma_0\subseteq\Gamma\) with the same derivation, and soundness
gives \(\Gamma_0\entails^{\ell}_{\mathcal C_L}A\). Conversely,
every world satisfying \(\Gamma\) satisfies its subset \(\Gamma_0\),
so local consequence from that subset implies local consequence from
\(\Gamma\).

Finally, suppose every finite subset of \(\Gamma\) is satisfiable.
If \(\Gamma\) were inconsistent, a finite subset would derive
\(\bottom\). Soundness would make that finite subset unsatisfiable,
contrary to the hypothesis. Thus \(\Gamma\) is consistent, and its
maximal extension supplies a satisfying world in \(\mathcal C_L\) by
the matching strong completeness theorem. This argument does not choose
one model for each finite subset; it uses only the asserted existence of
those models and constructs the final model syntactically.
\end{proof}
\end{annexitem}
